\chapter{Gouvernance et suivi}
\label{ch:gouvernance}

La gouvernance reste légère pour un projet académique mené en solo. Les rituels de revue sont en place. Les décisions sont tracées. Les risques sont suivis avec la même cadence que les livrables.

\section{Comitologie}
\label{sec:comitologie}

Quatre instances structurent le suivi du projet. Chacune a un rôle précis et une cadence fixe.

\bridgezebra
\begin{longtable}{p{4.5cm}p{7cm}p{3cm}}
  \toprule
  \rowcolor{bridgeblue!90}
  {\color{white}\bfseries Instance} & {\color{white}\bfseries Rôle} & {\color{white}\bfseries Cadence} \\
  \midrule
  \endfirsthead
  \toprule
  \rowcolor{bridgeblue!90}
  {\color{white}\bfseries Instance} & {\color{white}\bfseries Rôle} & {\color{white}\bfseries Cadence} \\
  \midrule
  \endhead
  \bottomrule
  \endfoot
  Auto-revue auteur & Vérification quotidienne du code et des livrables, journal de bord & Quotidienne \\
  Point d'avancement encadrant & Suivi des jalons, arbitrage des hypothèses, validation des choix techniques & Hebdomadaire \\
  Pré-jury blanc & Répétition orale de la démo et du pitch, ajustements de timing & Une fois à J6 \\
  Jury final & Soutenance et évaluation académique TS6 EPHEC Brussels & 2 mars 2027 \\
\end{longtable}

\section{Suivi des décisions}
\label{sec:suivi-decisions}

Les décisions structurantes sont consignées dans \texttt{PROJECT.md} (section décisions verrouillées). Les modifications de fond du CDC sont tracées dans le journal des versions de la page de garde du présent document. Toute décision révisée fait l'objet d'une nouvelle entrée datée. Le statut DM, la stack Python, le périmètre des 4 simulateurs et la philosophie 3 questions / 2 phases sont des décisions verrouillées non négociables.

\section{Suivi des risques}
\label{sec:suivi-risques}

Les cinq risques projet du chapitre~\ref{ch:roadmap-risques} font l'objet d'un point hebdomadaire avec l'encadrant. Les risques techniques du chapitre~\ref{ch:securite-conformite} sont évalués à chaque livraison de lot. Un risque qui passe en niveau élevé déclenche un arbitrage immédiat sur le périmètre du lot en cours.

\section{Communication}
\label{sec:communication}

L'audience principale du projet est Matthieu Oliviers, encadrant TS6, et le jury TS6 de l'EPHEC Brussels. L'audience secondaire est constituée des camarades TS6 lors du pitch. Aucune communication externe ou industrielle n'est prévue à ce stade. Les livrables ne sont diffusés qu'au jury et conservés à des fins de portfolio académique.
